% Orion --- Theory / Problem Formalization section (systems-paper style)
% ---------------------------------------------------------------------------
% Usage:
%   1) As an includable section, add to your preamble:
%        \usepackage{amsmath,amssymb,amsthm,booktabs}
%        \newtheorem{lemma}{Lemma}
%        \newtheorem{assumption}{Assumption}
%      then in the body:  % Orion --- Theory / Problem Formalization section (systems-paper style)
% ---------------------------------------------------------------------------
% Usage:
%   1) As an includable section, add to your preamble:
%        \usepackage{amsmath,amssymb,amsthm,booktabs}
%        \newtheorem{lemma}{Lemma}
%        \newtheorem{assumption}{Assumption}
%      then in the body:  % Orion --- Theory / Problem Formalization section (systems-paper style)
% ---------------------------------------------------------------------------
% Usage:
%   1) As an includable section, add to your preamble:
%        \usepackage{amsmath,amssymb,amsthm,booktabs}
%        \newtheorem{lemma}{Lemma}
%        \newtheorem{assumption}{Assumption}
%      then in the body:  % Orion --- Theory / Problem Formalization section (systems-paper style)
% ---------------------------------------------------------------------------
% Usage:
%   1) As an includable section, add to your preamble:
%        \usepackage{amsmath,amssymb,amsthm,booktabs}
%        \newtheorem{lemma}{Lemma}
%        \newtheorem{assumption}{Assumption}
%      then in the body:  \input{theory_section}
%   2) To compile this file standalone, uncomment the STANDALONE block below.
%
% Note: written in English per systems-venue convention (OSDI/NSDI/SOSP/VLDB/ATC).
% ---------------------------------------------------------------------------

% ------------------------- STANDALONE (uncomment) --------------------------
% \documentclass[11pt]{article}
% \usepackage[margin=1in]{geometry}
% \usepackage{amsmath,amssymb,amsthm,booktabs}
% \newtheorem{lemma}{Lemma}
% \newtheorem{assumption}{Assumption}
% \newtheorem{definition}{Definition}
% \begin{document}
% ---------------------------------------------------------------------------

\section{Sharding as Balanced Fan-out Minimization}
\label{sec:theory}

\paragraph{The sharding problem (mechanism-independent).}
We first state the problem without reference to any index structure. Let
$X=\{x_1,\dots,x_N\}$ be the dataset, $q\sim Q$ a query drawn from the workload,
and $\mathrm{NN}_k(q)\subseteq X$ its true $k$ nearest neighbors under the target
metric; $R$ denotes the target recall, i.e.\ we require Recall@$k\ge R$. A
sharding is an assignment $\sigma:X\to\{1,\dots,P\}$; \emph{replication}, which
lets a point occupy several shards, is a relaxation we treat at the end of this
section. To return the true neighbors, a
query must touch every shard that holds one of them, so its \emph{required
fan-out} is
\[
\mathrm{fanout}^\star_\sigma(q)=\bigl|\{\sigma(x):x\in\mathrm{NN}_k(q)\}\bigr|.
\]
This quantity depends only on the data, the workload, and the assignment---not on
how any system routes. At fixed recall the per-query server work is
$\sum_{s\,\text{probed}}\Theta(\mathrm{ef}_s\log|\text{shard}|)$, so with one shard
per core the throughput at saturation scales as
$P/\bigl(\mathbb{E}_q[\mathrm{fanout}]\cdot\rho\cdot\bar c\bigr)$, where $\rho$ is
the \emph{load skew} (busiest-over-mean per-shard query load) and $\bar c$ the
mean per-shard search cost (\S\ref{sec:motivation}, claim~C1). Fan-out is the
first-order term---it fixes the aggregate work---so we take it as the objective;
the load skew $\rho$ and the per-shard search strengths $\mathrm{ef}_s$ are
second-order factors that the same mechanism also controls (developed below and
measured in \S\ref{sec:eval}). The first-order sharding problem is
\begin{equation}
\label{eq:objective}
\min_{\sigma}\ \mathbb{E}_{q}\!\left[\mathrm{fanout}^\star_\sigma(q)\right]
\quad\text{s.t.}\quad
(1-\varepsilon)\tfrac{N}{P}\le |\sigma^{-1}(i)| \le (1+\varepsilon)\tfrac{N}{P}\ \ \forall i .
\end{equation}
The band constrains shard \emph{size} $|\sigma^{-1}(i)|$, a tractable surrogate
for the per-shard query \emph{load} that actually caps throughput; on a graph
index the two diverge---equal-sized shards can draw unequal traffic---so we
enforce size balance offline and report the residual load skew $\rho$ as a
separate diagnostic rather than fold it into~\eqref{eq:objective}.
\eqref{eq:objective} is an instance of a classical combinatorial problem, which
lets us reuse its known hardness and algorithms. A query's neighbor set
$\mathrm{NN}_k(q)$ ties together up to $k$ points that must be retrieved together,
a relation an ordinary (two-endpoint) edge cannot express. We therefore model the
workload as a \emph{hypergraph} $H=(X,\{e_q\}_q)$ whose vertices are the data
points and whose \emph{hyperedges} $e_q=\mathrm{NN}_k(q)$ each span all neighbors
of one query, with weight $c(e)=\Pr_q[\mathrm{NN}_k(q)=e]$. For an assignment
$\sigma$, let $\lambda(e)=|\{\sigma(x):x\in e\}|$ be the number of distinct shards
a hyperedge spans; by construction this is exactly the query's fan-out,
$\lambda(e_q)=\mathrm{fanout}^\star_\sigma(q)$. Hypergraph partitioning
conventionally minimizes the \emph{connectivity} (or ``$\lambda\!-\!1$'') metric
$\sum_e c(e)\,(\lambda(e)-1)$, i.e.\ how many \emph{extra} shards each hyperedge
reaches beyond its first---the standard measure of cross-partition communication.
Since $\sum_e c(e)=1$, this metric equals
$\mathbb{E}_q[\mathrm{fanout}^\star_\sigma(q)]-1$, so it has the \emph{same
minimizer} as~\eqref{eq:objective}: minimizing expected fan-out \emph{is}
minimizing the $\lambda\!-\!1$ objective, and the balance band is the standard
$(1{+}\varepsilon)$ balance constraint. Hence~\eqref{eq:objective} is a
$(P,1{+}\varepsilon)$-balanced minimum-connectivity hypergraph partitioning
instance ($P$ parts, within a $(1{+}\varepsilon)$ balance band)---the problem
solved by tools such as hMETIS and KaHyPar~\cite{metis,kahypar}. (For a target
$R<1$, $\mathrm{NN}_k(q)$ is replaced by a recall-$R$ sufficient subset, which can
only lower fan-out; the full-neighbor version above is the conservative
reference.) Even for fixed
$P$ with the relaxed balance constraint this problem is NP-hard, and
inapproximable under strictly equal parts~\cite{andreev2006}, which is why Orion
optimizes a computable surrogate rather than~\eqref{eq:objective} itself, as the
next paragraphs develop.

\paragraph{Why the problem cannot be optimized directly.}
Two gaps prevent solving~\eqref{eq:objective} as stated. \emph{Offline}, the
hyperedges $\mathrm{NN}_k(q)$ are query-specific and the workload is not known in
advance; materializing ground-truth neighbors for a representative query sample
at scale is expensive, so the partitioner needs a computable, query-independent
proxy for which points are \emph{co-relevant}. \emph{Online}, $\mathrm{NN}_k(q)$
is precisely the unknown the query is trying to find, so a router must estimate
which shards hold the neighbors without an exhaustive scan. Any practical system
must supply a mechanism for both; the navigation graph is Orion's, and it is a
component of the \emph{solution}, not a premise of the problem.

\paragraph{Orion's navigation-graph instantiation.}
Orion introduces a single frozen \emph{upper navigation graph} $G=(V,E,w)$, where
$V$ are upper navigation points and $E$ is its directed adjacency. We define the
\emph{navigation traffic}
\begin{equation}
\label{eq:traffic}
w(u,v)=\mathbb{E}_{q}\bigl[\,\#\{\text{times $q$'s traversal explores }(u,v)\}\,\bigr],
\end{equation}
the expected number of times, \emph{per query}, that the frozen navigator
evaluates $v$ as a neighbor of an expanded node $u$. This is exactly the
search-induced evidence Orion already materializes as attachment votes; it is not
raw geometric distance. Its working hypothesis (\emph{search-induced topology}) is
that $w$ is a computable proxy for co-relevance: points reached together during
navigation tend to fall in a common query's neighbor set. This closes both gaps
with one artifact.

\emph{Offline}, Orion replaces the intractable hypergraph
objective~\eqref{eq:objective} by a tractable \emph{graph} surrogate on $G$.
For an upper membership $\pi:V\to\{1,\dots,P\}$, the
\emph{navigation-weighted cut}
$\sum_{(u,v):\pi(u)\neq\pi(v)}w(u,v)$ sums the traffic on edges whose endpoints
lie in different shards---intuitively, how often navigation would have to leave
one shard for another. Minimizing it under the same balance band is the standard
approach taken by multilevel partitioners~\cite{metis,kahypar}, and Orion
implements it in three classical stages:
\begin{enumerate}
\item \emph{Initialization.} A geometric $k$-means clustering of the upper
  vectors into $P$ shards yields a cheap first assignment that ignores graph
  structure but is already roughly balanced.
\item \emph{Move-based refinement.} Following the Kernighan--Lin /
  Fiduccia--Mattheyses family~\cite{kl1970,fm1982}, Orion repeatedly considers
  moving a single upper point from its current shard to another. A move's
  \emph{gain} is the change in weighted cut (plus a topology vote from
  navigation evidence); a move is accepted only if its gain is nonnegative (or
  repairs a balance violation) and the target shard still has spare capacity.
  Sequential, capacity-aware acceptance avoids the classic collapse where many
  nodes race into the same underfull shard in one synchronous round.
\item \emph{Balance repair.} If some shards still violate the
  $(1{\pm}\varepsilon)$ band after refinement, Orion solves a small
  \emph{transportation} (min-cost flow) problem: overloaded shards are sources,
  underloaded shards are sinks, and admissible arcs are those justified by
  navigation evidence. Shipping one unit of flow along an arc corresponds to
  reassigning one point (or a group of points); the optimum restores balance at
  minimal disruption to the cut.
\end{enumerate}
All three stages read only the frozen upper graph, so the offline cost stays
near-linear and the result is checksum-reproducible.

\emph{Online}, serving is two-stage: a global traversal of $G$ selects shards and
per-shard entry points $\mathrm{EP}(q)\subseteq\mathrm{Vis}(q)$ (the estimator of
the neighbor-bearing shards), after which each selected shard runs an
\emph{independent} local search whose \emph{strength} scales with the same
evidence: a shard receiving more of $q$'s entry points is searched at a larger
$\mathrm{ef}_s=\mathrm{ef}_0+\gamma\,|\mathrm{EP}(q)\cap\pi^{-1}(s)|$. The
navigation signal thus governs both \emph{which} shards are probed (the fan-out)
and \emph{how hard} each is searched (the effort budget $\sum_s\mathrm{ef}_s$),
concentrating work on the shards navigation implicates. The \emph{realized}
fan-out is
$\mathrm{fanout}^\pi(q)=|\{\pi(v):v\in\mathrm{EP}(q)\}|$; base-layer placement
follows $\pi$ under the same balance band.

\paragraph{A mechanism-side bound and a recall assumption.}
Two claims connect Orion's mechanism back to the problem. First, the offline
objective it actually optimizes---the navigation-weighted cut---upper-bounds the
\emph{realized} online fan-out (Lemma~\ref{lem:fanout}). Second, minimizing that
realized fan-out is meaningful only if the router still returns the neighbors,
which we state as an explicit, empirically testable assumption
(Assumption~\ref{ass:recall}) rather than fold into the definitions.

\begin{lemma}[Realized fan-out is dominated by the weighted cut]
\label{lem:fanout}
Let $\mathrm{cut}(\pi)=\{(u,v)\in E:\pi(u)\neq\pi(v)\}$ and let $w$ be the
navigation traffic of~\eqref{eq:traffic}. Then
\[
\mathbb{E}_{q}\!\left[\mathrm{fanout}^\pi(q)\right]
\ \le\ 1+\!\!\sum_{(u,v)\in\mathrm{cut}(\pi)}\!\! w(u,v).
\]
\end{lemma}

\begin{proof}
Fix a query $q$, let $\mathrm{Vis}(q)$ be the set of nodes its upper traversal
expands, and let $T(q)\subseteq E$ be the set of edges the traversal explores.
Best-first search discovers a node only by exploring an edge out of an
already-expanded node, so $(\mathrm{Vis}(q),T(q))$ is \emph{connected}: it
contains a spanning tree rooted at the entry node. We stress that we do not
assume the expansion order forms a walk in $G$---consecutive expanded nodes need
not be adjacent, since candidates are popped from a priority queue.

Let $m=\bigl|\{\pi(v):v\in\mathrm{Vis}(q)\}\bigr|$ be the number of shards the
visited set spans, and contract each shard's portion of $\mathrm{Vis}(q)$ into a
single supernode. The contracted multigraph has $m$ supernodes and inherits
connectivity, hence carries at least $m-1$ edges; each of them is an explored
edge whose endpoints lie in different shards. Therefore
\[
m\ \le\ 1+\bigl|\{(u,v)\in T(q):\pi(u)\neq\pi(v)\}\bigr|.
\]
Entry points are selected only among expanded nodes, so
$\mathrm{EP}(q)\subseteq\mathrm{Vis}(q)$ and $\mathrm{fanout}^\pi(q)\le m$. Taking
expectations over $q$ and applying~\eqref{eq:traffic} edgewise,
$\mathbb{E}_q\bigl|\{(u,v)\in T(q):\pi(u)\neq\pi(v)\}\bigr|
=\sum_{(u,v)\in\mathrm{cut}(\pi)}w(u,v)$, which yields the claim.
\end{proof}

The bound is loose in two safe directions: fan-out counts only the shards of the
selected entry points, a subset of the visited shards, and a traversal that
repeatedly explores edges between the same two shards pays each time in the cut
but only once in fan-out. Thus the navigation-weighted cut is a
\emph{conservative} upper bound on realized fan-out, so optimizing it never
underestimates online cost.

A natural refinement of the same idea normalizes each shard's cut by its
\emph{volume} (total incident traffic)
$\mathrm{vol}(V_i)=\sum_{u\in V_i}\sum_{v}w(u,v)$, yielding the
\emph{conductance} (or normalized-cut) of the partition:
\[
\phi(\pi)=\sum_{i=1}^{P}\frac{\sum_{(u,v):\pi(u)=i,\pi(v)\neq i}w(u,v)}{\mathrm{vol}(V_i)}.
\]
Each summand is the probability that a random navigation step taken inside shard
$i$ immediately leaves it, and $\phi$ sums these escape probabilities over shards
(the normalized cut of Shi and Malik~\cite{shimalik2000}). Low conductance
therefore means every
shard is a \emph{metastable} region of the upper traversal---a set navigation
enters and then tends to stay in. Under this reading, the search-induced-topology hypothesis
simply says: shards that are metastable for navigation are the ones whose
interior points are co-relevant for the same queries, and are therefore the
right blocks for~\eqref{eq:objective}.

\begin{assumption}[Navigation preserves recall]
\label{ass:recall}
For target recall $R$, with high probability over $q\sim Q$ the shards selected by
the upper traversal, $\{\pi(v):v\in\mathrm{EP}(q)\}$, contain a recall-$R$
sufficient subset of $\mathrm{NN}_k(q)$.
\end{assumption}

Under Assumption~\ref{ass:recall}, Orion's router attains the target recall, so
its realized fan-out $\mathrm{fanout}^\pi$ is an \emph{achievable} cost for the
sharding problem, and by Lemma~\ref{lem:fanout} that cost is controlled by the
offline cut Orion minimizes. The required fan-out
$\mathrm{fanout}^\star_\sigma$ remains the mechanism-free reference against which
we report the residual gap; neither direction of comparison is automatic, since a
router targeting $R<1$ may deliberately skip a shard that holds a true neighbor.
The navigation graph thus links the mechanism-independent
objective to a computable surrogate \emph{without} being assumed in the problem
statement; Assumption~\ref{ass:recall} is what \S\ref{sec:eval} measures, and it
can fail---in which case the surrogate is simply a worse proxy, not a redefinition
of the problem.

\paragraph{From theory to mechanism.}
Table~\ref{tab:theory-map} maps each quantity in the formalization to the
concrete Orion mechanism and to the metric we report, so the design and the
evaluation can be read against the same objects.

One further design choice maps to a standard graph-partitioning relaxation.
So far $\sigma$ assigns each point to a single shard (\emph{edge-cut} style:
the cut is paid by edges that cross the boundary). Orion's
\emph{multi-assignment} instead lets a point occupy several shards, which is
the \emph{vertex-cut} (or overlapping) relaxation: a point that sits on a
search-relevant boundary is \emph{replicated} into the shards that navigation
evidence nominates. Replication lowers $\lambda$ for the hyperedges that
contain that point---the query no longer needs an extra hop to reach a copy---at
the cost of storing more physical copies. Formally, replication generalizes
$\sigma$ to $X\to 2^{\{1,\dots,P\}}$, the required fan-out becomes the smallest
number of shards \emph{covering} $\mathrm{NN}_k(q)$, and shard load is measured in
physical copies: with $M=\sum_{x\in X}|\sigma(x)|$ total copies, the balance band
in~\eqref{eq:objective} becomes $(1\pm\varepsilon)M/P$. The expansion ratio $M/N$
is the price of this relaxation, and is reported alongside fan-out so the
trade-off is explicit.

\begin{table}[t]
\centering
\small
\begin{tabular}{@{}lll@{}}
\toprule
Theoretical quantity & Orion mechanism & Reported metric \\
\midrule
Weighted cut / conductance & CCNB offline objective & cross-shard traffic \\
$\mathbb{E}[\mathrm{fanout}^\pi]$ (Lem.~\ref{lem:fanout}) & selective routing & visited logical shards \\
Balance band $(1{\pm}\varepsilon)$ & fixed-$P$ balancing & max/mean, CV \\
Load skew $\rho$ (2nd order) & size balance as proxy & max/mean query load \\
Per-shard strength $\mathrm{ef}_s$ (2nd order) & effort-scaled local search & within-shard work \\
Vertex-cut replication & multi-assignment & expansion ratio \\
\bottomrule
\end{tabular}
\caption{Correspondence between the partitioning formalization
(\S\ref{sec:theory}) and Orion's system mechanisms and measured quantities.}
\label{tab:theory-map}
\end{table}

We do not claim an approximation guarantee; instead, \S\ref{sec:eval}
empirically validates Lemma~\ref{lem:fanout} by correlating the offline
navigation-weighted cut (and conductance) with measured fan-out and end-to-end
QPS across $P\in\{4,8,16,32\}$, and tests Assumption~\ref{ass:recall} by checking
that the selected shards cover the true neighbors at the target recall. It also
quantifies the two second-order factors: the residual load skew $\rho$ that
survives size balancing, and the within-shard work that the effort-scaled
$\mathrm{ef}_s$ saves against a uniform-$\mathrm{ef}$ baseline at matched recall.

% ------------------------- STANDALONE (uncomment) --------------------------
% \end{document}
% ---------------------------------------------------------------------------

%   2) To compile this file standalone, uncomment the STANDALONE block below.
%
% Note: written in English per systems-venue convention (OSDI/NSDI/SOSP/VLDB/ATC).
% ---------------------------------------------------------------------------

% ------------------------- STANDALONE (uncomment) --------------------------
% \documentclass[11pt]{article}
% \usepackage[margin=1in]{geometry}
% \usepackage{amsmath,amssymb,amsthm,booktabs}
% \newtheorem{lemma}{Lemma}
% \newtheorem{assumption}{Assumption}
% \newtheorem{definition}{Definition}
% \begin{document}
% ---------------------------------------------------------------------------

\section{Sharding as Balanced Fan-out Minimization}
\label{sec:theory}

\paragraph{The sharding problem (mechanism-independent).}
We first state the problem without reference to any index structure. Let
$X=\{x_1,\dots,x_N\}$ be the dataset, $q\sim Q$ a query drawn from the workload,
and $\mathrm{NN}_k(q)\subseteq X$ its true $k$ nearest neighbors under the target
metric; $R$ denotes the target recall, i.e.\ we require Recall@$k\ge R$. A
sharding is an assignment $\sigma:X\to\{1,\dots,P\}$; \emph{replication}, which
lets a point occupy several shards, is a relaxation we treat at the end of this
section. To return the true neighbors, a
query must touch every shard that holds one of them, so its \emph{required
fan-out} is
\[
\mathrm{fanout}^\star_\sigma(q)=\bigl|\{\sigma(x):x\in\mathrm{NN}_k(q)\}\bigr|.
\]
This quantity depends only on the data, the workload, and the assignment---not on
how any system routes. At fixed recall the per-query server work is
$\sum_{s\,\text{probed}}\Theta(\mathrm{ef}_s\log|\text{shard}|)$, so with one shard
per core the throughput at saturation scales as
$P/\bigl(\mathbb{E}_q[\mathrm{fanout}]\cdot\rho\cdot\bar c\bigr)$, where $\rho$ is
the \emph{load skew} (busiest-over-mean per-shard query load) and $\bar c$ the
mean per-shard search cost (\S\ref{sec:motivation}, claim~C1). Fan-out is the
first-order term---it fixes the aggregate work---so we take it as the objective;
the load skew $\rho$ and the per-shard search strengths $\mathrm{ef}_s$ are
second-order factors that the same mechanism also controls (developed below and
measured in \S\ref{sec:eval}). The first-order sharding problem is
\begin{equation}
\label{eq:objective}
\min_{\sigma}\ \mathbb{E}_{q}\!\left[\mathrm{fanout}^\star_\sigma(q)\right]
\quad\text{s.t.}\quad
(1-\varepsilon)\tfrac{N}{P}\le |\sigma^{-1}(i)| \le (1+\varepsilon)\tfrac{N}{P}\ \ \forall i .
\end{equation}
The band constrains shard \emph{size} $|\sigma^{-1}(i)|$, a tractable surrogate
for the per-shard query \emph{load} that actually caps throughput; on a graph
index the two diverge---equal-sized shards can draw unequal traffic---so we
enforce size balance offline and report the residual load skew $\rho$ as a
separate diagnostic rather than fold it into~\eqref{eq:objective}.
\eqref{eq:objective} is an instance of a classical combinatorial problem, which
lets us reuse its known hardness and algorithms. A query's neighbor set
$\mathrm{NN}_k(q)$ ties together up to $k$ points that must be retrieved together,
a relation an ordinary (two-endpoint) edge cannot express. We therefore model the
workload as a \emph{hypergraph} $H=(X,\{e_q\}_q)$ whose vertices are the data
points and whose \emph{hyperedges} $e_q=\mathrm{NN}_k(q)$ each span all neighbors
of one query, with weight $c(e)=\Pr_q[\mathrm{NN}_k(q)=e]$. For an assignment
$\sigma$, let $\lambda(e)=|\{\sigma(x):x\in e\}|$ be the number of distinct shards
a hyperedge spans; by construction this is exactly the query's fan-out,
$\lambda(e_q)=\mathrm{fanout}^\star_\sigma(q)$. Hypergraph partitioning
conventionally minimizes the \emph{connectivity} (or ``$\lambda\!-\!1$'') metric
$\sum_e c(e)\,(\lambda(e)-1)$, i.e.\ how many \emph{extra} shards each hyperedge
reaches beyond its first---the standard measure of cross-partition communication.
Since $\sum_e c(e)=1$, this metric equals
$\mathbb{E}_q[\mathrm{fanout}^\star_\sigma(q)]-1$, so it has the \emph{same
minimizer} as~\eqref{eq:objective}: minimizing expected fan-out \emph{is}
minimizing the $\lambda\!-\!1$ objective, and the balance band is the standard
$(1{+}\varepsilon)$ balance constraint. Hence~\eqref{eq:objective} is a
$(P,1{+}\varepsilon)$-balanced minimum-connectivity hypergraph partitioning
instance ($P$ parts, within a $(1{+}\varepsilon)$ balance band)---the problem
solved by tools such as hMETIS and KaHyPar~\cite{metis,kahypar}. (For a target
$R<1$, $\mathrm{NN}_k(q)$ is replaced by a recall-$R$ sufficient subset, which can
only lower fan-out; the full-neighbor version above is the conservative
reference.) Even for fixed
$P$ with the relaxed balance constraint this problem is NP-hard, and
inapproximable under strictly equal parts~\cite{andreev2006}, which is why Orion
optimizes a computable surrogate rather than~\eqref{eq:objective} itself, as the
next paragraphs develop.

\paragraph{Why the problem cannot be optimized directly.}
Two gaps prevent solving~\eqref{eq:objective} as stated. \emph{Offline}, the
hyperedges $\mathrm{NN}_k(q)$ are query-specific and the workload is not known in
advance; materializing ground-truth neighbors for a representative query sample
at scale is expensive, so the partitioner needs a computable, query-independent
proxy for which points are \emph{co-relevant}. \emph{Online}, $\mathrm{NN}_k(q)$
is precisely the unknown the query is trying to find, so a router must estimate
which shards hold the neighbors without an exhaustive scan. Any practical system
must supply a mechanism for both; the navigation graph is Orion's, and it is a
component of the \emph{solution}, not a premise of the problem.

\paragraph{Orion's navigation-graph instantiation.}
Orion introduces a single frozen \emph{upper navigation graph} $G=(V,E,w)$, where
$V$ are upper navigation points and $E$ is its directed adjacency. We define the
\emph{navigation traffic}
\begin{equation}
\label{eq:traffic}
w(u,v)=\mathbb{E}_{q}\bigl[\,\#\{\text{times $q$'s traversal explores }(u,v)\}\,\bigr],
\end{equation}
the expected number of times, \emph{per query}, that the frozen navigator
evaluates $v$ as a neighbor of an expanded node $u$. This is exactly the
search-induced evidence Orion already materializes as attachment votes; it is not
raw geometric distance. Its working hypothesis (\emph{search-induced topology}) is
that $w$ is a computable proxy for co-relevance: points reached together during
navigation tend to fall in a common query's neighbor set. This closes both gaps
with one artifact.

\emph{Offline}, Orion replaces the intractable hypergraph
objective~\eqref{eq:objective} by a tractable \emph{graph} surrogate on $G$.
For an upper membership $\pi:V\to\{1,\dots,P\}$, the
\emph{navigation-weighted cut}
$\sum_{(u,v):\pi(u)\neq\pi(v)}w(u,v)$ sums the traffic on edges whose endpoints
lie in different shards---intuitively, how often navigation would have to leave
one shard for another. Minimizing it under the same balance band is the standard
approach taken by multilevel partitioners~\cite{metis,kahypar}, and Orion
implements it in three classical stages:
\begin{enumerate}
\item \emph{Initialization.} A geometric $k$-means clustering of the upper
  vectors into $P$ shards yields a cheap first assignment that ignores graph
  structure but is already roughly balanced.
\item \emph{Move-based refinement.} Following the Kernighan--Lin /
  Fiduccia--Mattheyses family~\cite{kl1970,fm1982}, Orion repeatedly considers
  moving a single upper point from its current shard to another. A move's
  \emph{gain} is the change in weighted cut (plus a topology vote from
  navigation evidence); a move is accepted only if its gain is nonnegative (or
  repairs a balance violation) and the target shard still has spare capacity.
  Sequential, capacity-aware acceptance avoids the classic collapse where many
  nodes race into the same underfull shard in one synchronous round.
\item \emph{Balance repair.} If some shards still violate the
  $(1{\pm}\varepsilon)$ band after refinement, Orion solves a small
  \emph{transportation} (min-cost flow) problem: overloaded shards are sources,
  underloaded shards are sinks, and admissible arcs are those justified by
  navigation evidence. Shipping one unit of flow along an arc corresponds to
  reassigning one point (or a group of points); the optimum restores balance at
  minimal disruption to the cut.
\end{enumerate}
All three stages read only the frozen upper graph, so the offline cost stays
near-linear and the result is checksum-reproducible.

\emph{Online}, serving is two-stage: a global traversal of $G$ selects shards and
per-shard entry points $\mathrm{EP}(q)\subseteq\mathrm{Vis}(q)$ (the estimator of
the neighbor-bearing shards), after which each selected shard runs an
\emph{independent} local search whose \emph{strength} scales with the same
evidence: a shard receiving more of $q$'s entry points is searched at a larger
$\mathrm{ef}_s=\mathrm{ef}_0+\gamma\,|\mathrm{EP}(q)\cap\pi^{-1}(s)|$. The
navigation signal thus governs both \emph{which} shards are probed (the fan-out)
and \emph{how hard} each is searched (the effort budget $\sum_s\mathrm{ef}_s$),
concentrating work on the shards navigation implicates. The \emph{realized}
fan-out is
$\mathrm{fanout}^\pi(q)=|\{\pi(v):v\in\mathrm{EP}(q)\}|$; base-layer placement
follows $\pi$ under the same balance band.

\paragraph{A mechanism-side bound and a recall assumption.}
Two claims connect Orion's mechanism back to the problem. First, the offline
objective it actually optimizes---the navigation-weighted cut---upper-bounds the
\emph{realized} online fan-out (Lemma~\ref{lem:fanout}). Second, minimizing that
realized fan-out is meaningful only if the router still returns the neighbors,
which we state as an explicit, empirically testable assumption
(Assumption~\ref{ass:recall}) rather than fold into the definitions.

\begin{lemma}[Realized fan-out is dominated by the weighted cut]
\label{lem:fanout}
Let $\mathrm{cut}(\pi)=\{(u,v)\in E:\pi(u)\neq\pi(v)\}$ and let $w$ be the
navigation traffic of~\eqref{eq:traffic}. Then
\[
\mathbb{E}_{q}\!\left[\mathrm{fanout}^\pi(q)\right]
\ \le\ 1+\!\!\sum_{(u,v)\in\mathrm{cut}(\pi)}\!\! w(u,v).
\]
\end{lemma}

\begin{proof}
Fix a query $q$, let $\mathrm{Vis}(q)$ be the set of nodes its upper traversal
expands, and let $T(q)\subseteq E$ be the set of edges the traversal explores.
Best-first search discovers a node only by exploring an edge out of an
already-expanded node, so $(\mathrm{Vis}(q),T(q))$ is \emph{connected}: it
contains a spanning tree rooted at the entry node. We stress that we do not
assume the expansion order forms a walk in $G$---consecutive expanded nodes need
not be adjacent, since candidates are popped from a priority queue.

Let $m=\bigl|\{\pi(v):v\in\mathrm{Vis}(q)\}\bigr|$ be the number of shards the
visited set spans, and contract each shard's portion of $\mathrm{Vis}(q)$ into a
single supernode. The contracted multigraph has $m$ supernodes and inherits
connectivity, hence carries at least $m-1$ edges; each of them is an explored
edge whose endpoints lie in different shards. Therefore
\[
m\ \le\ 1+\bigl|\{(u,v)\in T(q):\pi(u)\neq\pi(v)\}\bigr|.
\]
Entry points are selected only among expanded nodes, so
$\mathrm{EP}(q)\subseteq\mathrm{Vis}(q)$ and $\mathrm{fanout}^\pi(q)\le m$. Taking
expectations over $q$ and applying~\eqref{eq:traffic} edgewise,
$\mathbb{E}_q\bigl|\{(u,v)\in T(q):\pi(u)\neq\pi(v)\}\bigr|
=\sum_{(u,v)\in\mathrm{cut}(\pi)}w(u,v)$, which yields the claim.
\end{proof}

The bound is loose in two safe directions: fan-out counts only the shards of the
selected entry points, a subset of the visited shards, and a traversal that
repeatedly explores edges between the same two shards pays each time in the cut
but only once in fan-out. Thus the navigation-weighted cut is a
\emph{conservative} upper bound on realized fan-out, so optimizing it never
underestimates online cost.

A natural refinement of the same idea normalizes each shard's cut by its
\emph{volume} (total incident traffic)
$\mathrm{vol}(V_i)=\sum_{u\in V_i}\sum_{v}w(u,v)$, yielding the
\emph{conductance} (or normalized-cut) of the partition:
\[
\phi(\pi)=\sum_{i=1}^{P}\frac{\sum_{(u,v):\pi(u)=i,\pi(v)\neq i}w(u,v)}{\mathrm{vol}(V_i)}.
\]
Each summand is the probability that a random navigation step taken inside shard
$i$ immediately leaves it, and $\phi$ sums these escape probabilities over shards
(the normalized cut of Shi and Malik~\cite{shimalik2000}). Low conductance
therefore means every
shard is a \emph{metastable} region of the upper traversal---a set navigation
enters and then tends to stay in. Under this reading, the search-induced-topology hypothesis
simply says: shards that are metastable for navigation are the ones whose
interior points are co-relevant for the same queries, and are therefore the
right blocks for~\eqref{eq:objective}.

\begin{assumption}[Navigation preserves recall]
\label{ass:recall}
For target recall $R$, with high probability over $q\sim Q$ the shards selected by
the upper traversal, $\{\pi(v):v\in\mathrm{EP}(q)\}$, contain a recall-$R$
sufficient subset of $\mathrm{NN}_k(q)$.
\end{assumption}

Under Assumption~\ref{ass:recall}, Orion's router attains the target recall, so
its realized fan-out $\mathrm{fanout}^\pi$ is an \emph{achievable} cost for the
sharding problem, and by Lemma~\ref{lem:fanout} that cost is controlled by the
offline cut Orion minimizes. The required fan-out
$\mathrm{fanout}^\star_\sigma$ remains the mechanism-free reference against which
we report the residual gap; neither direction of comparison is automatic, since a
router targeting $R<1$ may deliberately skip a shard that holds a true neighbor.
The navigation graph thus links the mechanism-independent
objective to a computable surrogate \emph{without} being assumed in the problem
statement; Assumption~\ref{ass:recall} is what \S\ref{sec:eval} measures, and it
can fail---in which case the surrogate is simply a worse proxy, not a redefinition
of the problem.

\paragraph{From theory to mechanism.}
Table~\ref{tab:theory-map} maps each quantity in the formalization to the
concrete Orion mechanism and to the metric we report, so the design and the
evaluation can be read against the same objects.

One further design choice maps to a standard graph-partitioning relaxation.
So far $\sigma$ assigns each point to a single shard (\emph{edge-cut} style:
the cut is paid by edges that cross the boundary). Orion's
\emph{multi-assignment} instead lets a point occupy several shards, which is
the \emph{vertex-cut} (or overlapping) relaxation: a point that sits on a
search-relevant boundary is \emph{replicated} into the shards that navigation
evidence nominates. Replication lowers $\lambda$ for the hyperedges that
contain that point---the query no longer needs an extra hop to reach a copy---at
the cost of storing more physical copies. Formally, replication generalizes
$\sigma$ to $X\to 2^{\{1,\dots,P\}}$, the required fan-out becomes the smallest
number of shards \emph{covering} $\mathrm{NN}_k(q)$, and shard load is measured in
physical copies: with $M=\sum_{x\in X}|\sigma(x)|$ total copies, the balance band
in~\eqref{eq:objective} becomes $(1\pm\varepsilon)M/P$. The expansion ratio $M/N$
is the price of this relaxation, and is reported alongside fan-out so the
trade-off is explicit.

\begin{table}[t]
\centering
\small
\begin{tabular}{@{}lll@{}}
\toprule
Theoretical quantity & Orion mechanism & Reported metric \\
\midrule
Weighted cut / conductance & CCNB offline objective & cross-shard traffic \\
$\mathbb{E}[\mathrm{fanout}^\pi]$ (Lem.~\ref{lem:fanout}) & selective routing & visited logical shards \\
Balance band $(1{\pm}\varepsilon)$ & fixed-$P$ balancing & max/mean, CV \\
Load skew $\rho$ (2nd order) & size balance as proxy & max/mean query load \\
Per-shard strength $\mathrm{ef}_s$ (2nd order) & effort-scaled local search & within-shard work \\
Vertex-cut replication & multi-assignment & expansion ratio \\
\bottomrule
\end{tabular}
\caption{Correspondence between the partitioning formalization
(\S\ref{sec:theory}) and Orion's system mechanisms and measured quantities.}
\label{tab:theory-map}
\end{table}

We do not claim an approximation guarantee; instead, \S\ref{sec:eval}
empirically validates Lemma~\ref{lem:fanout} by correlating the offline
navigation-weighted cut (and conductance) with measured fan-out and end-to-end
QPS across $P\in\{4,8,16,32\}$, and tests Assumption~\ref{ass:recall} by checking
that the selected shards cover the true neighbors at the target recall. It also
quantifies the two second-order factors: the residual load skew $\rho$ that
survives size balancing, and the within-shard work that the effort-scaled
$\mathrm{ef}_s$ saves against a uniform-$\mathrm{ef}$ baseline at matched recall.

% ------------------------- STANDALONE (uncomment) --------------------------
% \end{document}
% ---------------------------------------------------------------------------

%   2) To compile this file standalone, uncomment the STANDALONE block below.
%
% Note: written in English per systems-venue convention (OSDI/NSDI/SOSP/VLDB/ATC).
% ---------------------------------------------------------------------------

% ------------------------- STANDALONE (uncomment) --------------------------
% \documentclass[11pt]{article}
% \usepackage[margin=1in]{geometry}
% \usepackage{amsmath,amssymb,amsthm,booktabs}
% \newtheorem{lemma}{Lemma}
% \newtheorem{assumption}{Assumption}
% \newtheorem{definition}{Definition}
% \begin{document}
% ---------------------------------------------------------------------------

\section{Sharding as Balanced Fan-out Minimization}
\label{sec:theory}

\paragraph{The sharding problem (mechanism-independent).}
We first state the problem without reference to any index structure. Let
$X=\{x_1,\dots,x_N\}$ be the dataset, $q\sim Q$ a query drawn from the workload,
and $\mathrm{NN}_k(q)\subseteq X$ its true $k$ nearest neighbors under the target
metric; $R$ denotes the target recall, i.e.\ we require Recall@$k\ge R$. A
sharding is an assignment $\sigma:X\to\{1,\dots,P\}$; \emph{replication}, which
lets a point occupy several shards, is a relaxation we treat at the end of this
section. To return the true neighbors, a
query must touch every shard that holds one of them, so its \emph{required
fan-out} is
\[
\mathrm{fanout}^\star_\sigma(q)=\bigl|\{\sigma(x):x\in\mathrm{NN}_k(q)\}\bigr|.
\]
This quantity depends only on the data, the workload, and the assignment---not on
how any system routes. At fixed recall the per-query server work is
$\sum_{s\,\text{probed}}\Theta(\mathrm{ef}_s\log|\text{shard}|)$, so with one shard
per core the throughput at saturation scales as
$P/\bigl(\mathbb{E}_q[\mathrm{fanout}]\cdot\rho\cdot\bar c\bigr)$, where $\rho$ is
the \emph{load skew} (busiest-over-mean per-shard query load) and $\bar c$ the
mean per-shard search cost (\S\ref{sec:motivation}, claim~C1). Fan-out is the
first-order term---it fixes the aggregate work---so we take it as the objective;
the load skew $\rho$ and the per-shard search strengths $\mathrm{ef}_s$ are
second-order factors that the same mechanism also controls (developed below and
measured in \S\ref{sec:eval}). The first-order sharding problem is
\begin{equation}
\label{eq:objective}
\min_{\sigma}\ \mathbb{E}_{q}\!\left[\mathrm{fanout}^\star_\sigma(q)\right]
\quad\text{s.t.}\quad
(1-\varepsilon)\tfrac{N}{P}\le |\sigma^{-1}(i)| \le (1+\varepsilon)\tfrac{N}{P}\ \ \forall i .
\end{equation}
The band constrains shard \emph{size} $|\sigma^{-1}(i)|$, a tractable surrogate
for the per-shard query \emph{load} that actually caps throughput; on a graph
index the two diverge---equal-sized shards can draw unequal traffic---so we
enforce size balance offline and report the residual load skew $\rho$ as a
separate diagnostic rather than fold it into~\eqref{eq:objective}.
\eqref{eq:objective} is an instance of a classical combinatorial problem, which
lets us reuse its known hardness and algorithms. A query's neighbor set
$\mathrm{NN}_k(q)$ ties together up to $k$ points that must be retrieved together,
a relation an ordinary (two-endpoint) edge cannot express. We therefore model the
workload as a \emph{hypergraph} $H=(X,\{e_q\}_q)$ whose vertices are the data
points and whose \emph{hyperedges} $e_q=\mathrm{NN}_k(q)$ each span all neighbors
of one query, with weight $c(e)=\Pr_q[\mathrm{NN}_k(q)=e]$. For an assignment
$\sigma$, let $\lambda(e)=|\{\sigma(x):x\in e\}|$ be the number of distinct shards
a hyperedge spans; by construction this is exactly the query's fan-out,
$\lambda(e_q)=\mathrm{fanout}^\star_\sigma(q)$. Hypergraph partitioning
conventionally minimizes the \emph{connectivity} (or ``$\lambda\!-\!1$'') metric
$\sum_e c(e)\,(\lambda(e)-1)$, i.e.\ how many \emph{extra} shards each hyperedge
reaches beyond its first---the standard measure of cross-partition communication.
Since $\sum_e c(e)=1$, this metric equals
$\mathbb{E}_q[\mathrm{fanout}^\star_\sigma(q)]-1$, so it has the \emph{same
minimizer} as~\eqref{eq:objective}: minimizing expected fan-out \emph{is}
minimizing the $\lambda\!-\!1$ objective, and the balance band is the standard
$(1{+}\varepsilon)$ balance constraint. Hence~\eqref{eq:objective} is a
$(P,1{+}\varepsilon)$-balanced minimum-connectivity hypergraph partitioning
instance ($P$ parts, within a $(1{+}\varepsilon)$ balance band)---the problem
solved by tools such as hMETIS and KaHyPar~\cite{metis,kahypar}. (For a target
$R<1$, $\mathrm{NN}_k(q)$ is replaced by a recall-$R$ sufficient subset, which can
only lower fan-out; the full-neighbor version above is the conservative
reference.) Even for fixed
$P$ with the relaxed balance constraint this problem is NP-hard, and
inapproximable under strictly equal parts~\cite{andreev2006}, which is why Orion
optimizes a computable surrogate rather than~\eqref{eq:objective} itself, as the
next paragraphs develop.

\paragraph{Why the problem cannot be optimized directly.}
Two gaps prevent solving~\eqref{eq:objective} as stated. \emph{Offline}, the
hyperedges $\mathrm{NN}_k(q)$ are query-specific and the workload is not known in
advance; materializing ground-truth neighbors for a representative query sample
at scale is expensive, so the partitioner needs a computable, query-independent
proxy for which points are \emph{co-relevant}. \emph{Online}, $\mathrm{NN}_k(q)$
is precisely the unknown the query is trying to find, so a router must estimate
which shards hold the neighbors without an exhaustive scan. Any practical system
must supply a mechanism for both; the navigation graph is Orion's, and it is a
component of the \emph{solution}, not a premise of the problem.

\paragraph{Orion's navigation-graph instantiation.}
Orion introduces a single frozen \emph{upper navigation graph} $G=(V,E,w)$, where
$V$ are upper navigation points and $E$ is its directed adjacency. We define the
\emph{navigation traffic}
\begin{equation}
\label{eq:traffic}
w(u,v)=\mathbb{E}_{q}\bigl[\,\#\{\text{times $q$'s traversal explores }(u,v)\}\,\bigr],
\end{equation}
the expected number of times, \emph{per query}, that the frozen navigator
evaluates $v$ as a neighbor of an expanded node $u$. This is exactly the
search-induced evidence Orion already materializes as attachment votes; it is not
raw geometric distance. Its working hypothesis (\emph{search-induced topology}) is
that $w$ is a computable proxy for co-relevance: points reached together during
navigation tend to fall in a common query's neighbor set. This closes both gaps
with one artifact.

\emph{Offline}, Orion replaces the intractable hypergraph
objective~\eqref{eq:objective} by a tractable \emph{graph} surrogate on $G$.
For an upper membership $\pi:V\to\{1,\dots,P\}$, the
\emph{navigation-weighted cut}
$\sum_{(u,v):\pi(u)\neq\pi(v)}w(u,v)$ sums the traffic on edges whose endpoints
lie in different shards---intuitively, how often navigation would have to leave
one shard for another. Minimizing it under the same balance band is the standard
approach taken by multilevel partitioners~\cite{metis,kahypar}, and Orion
implements it in three classical stages:
\begin{enumerate}
\item \emph{Initialization.} A geometric $k$-means clustering of the upper
  vectors into $P$ shards yields a cheap first assignment that ignores graph
  structure but is already roughly balanced.
\item \emph{Move-based refinement.} Following the Kernighan--Lin /
  Fiduccia--Mattheyses family~\cite{kl1970,fm1982}, Orion repeatedly considers
  moving a single upper point from its current shard to another. A move's
  \emph{gain} is the change in weighted cut (plus a topology vote from
  navigation evidence); a move is accepted only if its gain is nonnegative (or
  repairs a balance violation) and the target shard still has spare capacity.
  Sequential, capacity-aware acceptance avoids the classic collapse where many
  nodes race into the same underfull shard in one synchronous round.
\item \emph{Balance repair.} If some shards still violate the
  $(1{\pm}\varepsilon)$ band after refinement, Orion solves a small
  \emph{transportation} (min-cost flow) problem: overloaded shards are sources,
  underloaded shards are sinks, and admissible arcs are those justified by
  navigation evidence. Shipping one unit of flow along an arc corresponds to
  reassigning one point (or a group of points); the optimum restores balance at
  minimal disruption to the cut.
\end{enumerate}
All three stages read only the frozen upper graph, so the offline cost stays
near-linear and the result is checksum-reproducible.

\emph{Online}, serving is two-stage: a global traversal of $G$ selects shards and
per-shard entry points $\mathrm{EP}(q)\subseteq\mathrm{Vis}(q)$ (the estimator of
the neighbor-bearing shards), after which each selected shard runs an
\emph{independent} local search whose \emph{strength} scales with the same
evidence: a shard receiving more of $q$'s entry points is searched at a larger
$\mathrm{ef}_s=\mathrm{ef}_0+\gamma\,|\mathrm{EP}(q)\cap\pi^{-1}(s)|$. The
navigation signal thus governs both \emph{which} shards are probed (the fan-out)
and \emph{how hard} each is searched (the effort budget $\sum_s\mathrm{ef}_s$),
concentrating work on the shards navigation implicates. The \emph{realized}
fan-out is
$\mathrm{fanout}^\pi(q)=|\{\pi(v):v\in\mathrm{EP}(q)\}|$; base-layer placement
follows $\pi$ under the same balance band.

\paragraph{A mechanism-side bound and a recall assumption.}
Two claims connect Orion's mechanism back to the problem. First, the offline
objective it actually optimizes---the navigation-weighted cut---upper-bounds the
\emph{realized} online fan-out (Lemma~\ref{lem:fanout}). Second, minimizing that
realized fan-out is meaningful only if the router still returns the neighbors,
which we state as an explicit, empirically testable assumption
(Assumption~\ref{ass:recall}) rather than fold into the definitions.

\begin{lemma}[Realized fan-out is dominated by the weighted cut]
\label{lem:fanout}
Let $\mathrm{cut}(\pi)=\{(u,v)\in E:\pi(u)\neq\pi(v)\}$ and let $w$ be the
navigation traffic of~\eqref{eq:traffic}. Then
\[
\mathbb{E}_{q}\!\left[\mathrm{fanout}^\pi(q)\right]
\ \le\ 1+\!\!\sum_{(u,v)\in\mathrm{cut}(\pi)}\!\! w(u,v).
\]
\end{lemma}

\begin{proof}
Fix a query $q$, let $\mathrm{Vis}(q)$ be the set of nodes its upper traversal
expands, and let $T(q)\subseteq E$ be the set of edges the traversal explores.
Best-first search discovers a node only by exploring an edge out of an
already-expanded node, so $(\mathrm{Vis}(q),T(q))$ is \emph{connected}: it
contains a spanning tree rooted at the entry node. We stress that we do not
assume the expansion order forms a walk in $G$---consecutive expanded nodes need
not be adjacent, since candidates are popped from a priority queue.

Let $m=\bigl|\{\pi(v):v\in\mathrm{Vis}(q)\}\bigr|$ be the number of shards the
visited set spans, and contract each shard's portion of $\mathrm{Vis}(q)$ into a
single supernode. The contracted multigraph has $m$ supernodes and inherits
connectivity, hence carries at least $m-1$ edges; each of them is an explored
edge whose endpoints lie in different shards. Therefore
\[
m\ \le\ 1+\bigl|\{(u,v)\in T(q):\pi(u)\neq\pi(v)\}\bigr|.
\]
Entry points are selected only among expanded nodes, so
$\mathrm{EP}(q)\subseteq\mathrm{Vis}(q)$ and $\mathrm{fanout}^\pi(q)\le m$. Taking
expectations over $q$ and applying~\eqref{eq:traffic} edgewise,
$\mathbb{E}_q\bigl|\{(u,v)\in T(q):\pi(u)\neq\pi(v)\}\bigr|
=\sum_{(u,v)\in\mathrm{cut}(\pi)}w(u,v)$, which yields the claim.
\end{proof}

The bound is loose in two safe directions: fan-out counts only the shards of the
selected entry points, a subset of the visited shards, and a traversal that
repeatedly explores edges between the same two shards pays each time in the cut
but only once in fan-out. Thus the navigation-weighted cut is a
\emph{conservative} upper bound on realized fan-out, so optimizing it never
underestimates online cost.

A natural refinement of the same idea normalizes each shard's cut by its
\emph{volume} (total incident traffic)
$\mathrm{vol}(V_i)=\sum_{u\in V_i}\sum_{v}w(u,v)$, yielding the
\emph{conductance} (or normalized-cut) of the partition:
\[
\phi(\pi)=\sum_{i=1}^{P}\frac{\sum_{(u,v):\pi(u)=i,\pi(v)\neq i}w(u,v)}{\mathrm{vol}(V_i)}.
\]
Each summand is the probability that a random navigation step taken inside shard
$i$ immediately leaves it, and $\phi$ sums these escape probabilities over shards
(the normalized cut of Shi and Malik~\cite{shimalik2000}). Low conductance
therefore means every
shard is a \emph{metastable} region of the upper traversal---a set navigation
enters and then tends to stay in. Under this reading, the search-induced-topology hypothesis
simply says: shards that are metastable for navigation are the ones whose
interior points are co-relevant for the same queries, and are therefore the
right blocks for~\eqref{eq:objective}.

\begin{assumption}[Navigation preserves recall]
\label{ass:recall}
For target recall $R$, with high probability over $q\sim Q$ the shards selected by
the upper traversal, $\{\pi(v):v\in\mathrm{EP}(q)\}$, contain a recall-$R$
sufficient subset of $\mathrm{NN}_k(q)$.
\end{assumption}

Under Assumption~\ref{ass:recall}, Orion's router attains the target recall, so
its realized fan-out $\mathrm{fanout}^\pi$ is an \emph{achievable} cost for the
sharding problem, and by Lemma~\ref{lem:fanout} that cost is controlled by the
offline cut Orion minimizes. The required fan-out
$\mathrm{fanout}^\star_\sigma$ remains the mechanism-free reference against which
we report the residual gap; neither direction of comparison is automatic, since a
router targeting $R<1$ may deliberately skip a shard that holds a true neighbor.
The navigation graph thus links the mechanism-independent
objective to a computable surrogate \emph{without} being assumed in the problem
statement; Assumption~\ref{ass:recall} is what \S\ref{sec:eval} measures, and it
can fail---in which case the surrogate is simply a worse proxy, not a redefinition
of the problem.

\paragraph{From theory to mechanism.}
Table~\ref{tab:theory-map} maps each quantity in the formalization to the
concrete Orion mechanism and to the metric we report, so the design and the
evaluation can be read against the same objects.

One further design choice maps to a standard graph-partitioning relaxation.
So far $\sigma$ assigns each point to a single shard (\emph{edge-cut} style:
the cut is paid by edges that cross the boundary). Orion's
\emph{multi-assignment} instead lets a point occupy several shards, which is
the \emph{vertex-cut} (or overlapping) relaxation: a point that sits on a
search-relevant boundary is \emph{replicated} into the shards that navigation
evidence nominates. Replication lowers $\lambda$ for the hyperedges that
contain that point---the query no longer needs an extra hop to reach a copy---at
the cost of storing more physical copies. Formally, replication generalizes
$\sigma$ to $X\to 2^{\{1,\dots,P\}}$, the required fan-out becomes the smallest
number of shards \emph{covering} $\mathrm{NN}_k(q)$, and shard load is measured in
physical copies: with $M=\sum_{x\in X}|\sigma(x)|$ total copies, the balance band
in~\eqref{eq:objective} becomes $(1\pm\varepsilon)M/P$. The expansion ratio $M/N$
is the price of this relaxation, and is reported alongside fan-out so the
trade-off is explicit.

\begin{table}[t]
\centering
\small
\begin{tabular}{@{}lll@{}}
\toprule
Theoretical quantity & Orion mechanism & Reported metric \\
\midrule
Weighted cut / conductance & CCNB offline objective & cross-shard traffic \\
$\mathbb{E}[\mathrm{fanout}^\pi]$ (Lem.~\ref{lem:fanout}) & selective routing & visited logical shards \\
Balance band $(1{\pm}\varepsilon)$ & fixed-$P$ balancing & max/mean, CV \\
Load skew $\rho$ (2nd order) & size balance as proxy & max/mean query load \\
Per-shard strength $\mathrm{ef}_s$ (2nd order) & effort-scaled local search & within-shard work \\
Vertex-cut replication & multi-assignment & expansion ratio \\
\bottomrule
\end{tabular}
\caption{Correspondence between the partitioning formalization
(\S\ref{sec:theory}) and Orion's system mechanisms and measured quantities.}
\label{tab:theory-map}
\end{table}

We do not claim an approximation guarantee; instead, \S\ref{sec:eval}
empirically validates Lemma~\ref{lem:fanout} by correlating the offline
navigation-weighted cut (and conductance) with measured fan-out and end-to-end
QPS across $P\in\{4,8,16,32\}$, and tests Assumption~\ref{ass:recall} by checking
that the selected shards cover the true neighbors at the target recall. It also
quantifies the two second-order factors: the residual load skew $\rho$ that
survives size balancing, and the within-shard work that the effort-scaled
$\mathrm{ef}_s$ saves against a uniform-$\mathrm{ef}$ baseline at matched recall.

% ------------------------- STANDALONE (uncomment) --------------------------
% \end{document}
% ---------------------------------------------------------------------------

%   2) To compile this file standalone, uncomment the STANDALONE block below.
%
% Note: written in English per systems-venue convention (OSDI/NSDI/SOSP/VLDB/ATC).
% ---------------------------------------------------------------------------

% ------------------------- STANDALONE (uncomment) --------------------------
% \documentclass[11pt]{article}
% \usepackage[margin=1in]{geometry}
% \usepackage{amsmath,amssymb,amsthm,booktabs}
% \newtheorem{lemma}{Lemma}
% \newtheorem{assumption}{Assumption}
% \newtheorem{definition}{Definition}
% \begin{document}
% ---------------------------------------------------------------------------

\section{Sharding as Balanced Fan-out Minimization}
\label{sec:theory}

\paragraph{The sharding problem (mechanism-independent).}
We first state the problem without reference to any index structure. Let
$X=\{x_1,\dots,x_N\}$ be the dataset, $q\sim Q$ a query drawn from the workload,
and $\mathrm{NN}_k(q)\subseteq X$ its true $k$ nearest neighbors under the target
metric; $R$ denotes the target recall, i.e.\ we require Recall@$k\ge R$. A
sharding is an assignment $\sigma:X\to\{1,\dots,P\}$; \emph{replication}, which
lets a point occupy several shards, is a relaxation we treat at the end of this
section. To return the true neighbors, a
query must touch every shard that holds one of them, so its \emph{required
fan-out} is
\[
\mathrm{fanout}^\star_\sigma(q)=\bigl|\{\sigma(x):x\in\mathrm{NN}_k(q)\}\bigr|.
\]
This quantity depends only on the data, the workload, and the assignment---not on
how any system routes. At fixed recall the per-query server work is
$\sum_{s\,\text{probed}}\Theta(\mathrm{ef}_s\log|\text{shard}|)$, so with one shard
per core the throughput at saturation scales as
$P/\bigl(\mathbb{E}_q[\mathrm{fanout}]\cdot\rho\cdot\bar c\bigr)$, where $\rho$ is
the \emph{load skew} (busiest-over-mean per-shard query load) and $\bar c$ the
mean per-shard search cost (\S\ref{sec:motivation}, claim~C1). Fan-out is the
first-order term---it fixes the aggregate work---so we take it as the objective;
the load skew $\rho$ and the per-shard search strengths $\mathrm{ef}_s$ are
second-order factors that the same mechanism also controls (developed below and
measured in \S\ref{sec:eval}). The first-order sharding problem is
\begin{equation}
\label{eq:objective}
\min_{\sigma}\ \mathbb{E}_{q}\!\left[\mathrm{fanout}^\star_\sigma(q)\right]
\quad\text{s.t.}\quad
(1-\varepsilon)\tfrac{N}{P}\le |\sigma^{-1}(i)| \le (1+\varepsilon)\tfrac{N}{P}\ \ \forall i .
\end{equation}
The band constrains shard \emph{size} $|\sigma^{-1}(i)|$, a tractable surrogate
for the per-shard query \emph{load} that actually caps throughput; on a graph
index the two diverge---equal-sized shards can draw unequal traffic---so we
enforce size balance offline and report the residual load skew $\rho$ as a
separate diagnostic rather than fold it into~\eqref{eq:objective}.
\eqref{eq:objective} is an instance of a classical combinatorial problem, which
lets us reuse its known hardness and algorithms. A query's neighbor set
$\mathrm{NN}_k(q)$ ties together up to $k$ points that must be retrieved together,
a relation an ordinary (two-endpoint) edge cannot express. We therefore model the
workload as a \emph{hypergraph} $H=(X,\{e_q\}_q)$ whose vertices are the data
points and whose \emph{hyperedges} $e_q=\mathrm{NN}_k(q)$ each span all neighbors
of one query, with weight $c(e)=\Pr_q[\mathrm{NN}_k(q)=e]$. For an assignment
$\sigma$, let $\lambda(e)=|\{\sigma(x):x\in e\}|$ be the number of distinct shards
a hyperedge spans; by construction this is exactly the query's fan-out,
$\lambda(e_q)=\mathrm{fanout}^\star_\sigma(q)$. Hypergraph partitioning
conventionally minimizes the \emph{connectivity} (or ``$\lambda\!-\!1$'') metric
$\sum_e c(e)\,(\lambda(e)-1)$, i.e.\ how many \emph{extra} shards each hyperedge
reaches beyond its first---the standard measure of cross-partition communication.
Since $\sum_e c(e)=1$, this metric equals
$\mathbb{E}_q[\mathrm{fanout}^\star_\sigma(q)]-1$, so it has the \emph{same
minimizer} as~\eqref{eq:objective}: minimizing expected fan-out \emph{is}
minimizing the $\lambda\!-\!1$ objective, and the balance band is the standard
$(1{+}\varepsilon)$ balance constraint. Hence~\eqref{eq:objective} is a
$(P,1{+}\varepsilon)$-balanced minimum-connectivity hypergraph partitioning
instance ($P$ parts, within a $(1{+}\varepsilon)$ balance band)---the problem
solved by tools such as hMETIS and KaHyPar~\cite{metis,kahypar}. (For a target
$R<1$, $\mathrm{NN}_k(q)$ is replaced by a recall-$R$ sufficient subset, which can
only lower fan-out; the full-neighbor version above is the conservative
reference.) Even for fixed
$P$ with the relaxed balance constraint this problem is NP-hard, and
inapproximable under strictly equal parts~\cite{andreev2006}, which is why Orion
optimizes a computable surrogate rather than~\eqref{eq:objective} itself, as the
next paragraphs develop.

\paragraph{Why the problem cannot be optimized directly.}
Two gaps prevent solving~\eqref{eq:objective} as stated. \emph{Offline}, the
hyperedges $\mathrm{NN}_k(q)$ are query-specific and the workload is not known in
advance; materializing ground-truth neighbors for a representative query sample
at scale is expensive, so the partitioner needs a computable, query-independent
proxy for which points are \emph{co-relevant}. \emph{Online}, $\mathrm{NN}_k(q)$
is precisely the unknown the query is trying to find, so a router must estimate
which shards hold the neighbors without an exhaustive scan. Any practical system
must supply a mechanism for both; the navigation graph is Orion's, and it is a
component of the \emph{solution}, not a premise of the problem.

\paragraph{Orion's navigation-graph instantiation.}
Orion introduces a single frozen \emph{upper navigation graph} $G=(V,E,w)$, where
$V$ are upper navigation points and $E$ is its directed adjacency. We define the
\emph{navigation traffic}
\begin{equation}
\label{eq:traffic}
w(u,v)=\mathbb{E}_{q}\bigl[\,\#\{\text{times $q$'s traversal explores }(u,v)\}\,\bigr],
\end{equation}
the expected number of times, \emph{per query}, that the frozen navigator
evaluates $v$ as a neighbor of an expanded node $u$. This is exactly the
search-induced evidence Orion already materializes as attachment votes; it is not
raw geometric distance. Its working hypothesis (\emph{search-induced topology}) is
that $w$ is a computable proxy for co-relevance: points reached together during
navigation tend to fall in a common query's neighbor set. This closes both gaps
with one artifact.

\emph{Offline}, Orion replaces the intractable hypergraph
objective~\eqref{eq:objective} by a tractable \emph{graph} surrogate on $G$.
For an upper membership $\pi:V\to\{1,\dots,P\}$, the
\emph{navigation-weighted cut}
$\sum_{(u,v):\pi(u)\neq\pi(v)}w(u,v)$ sums the traffic on edges whose endpoints
lie in different shards---intuitively, how often navigation would have to leave
one shard for another. Minimizing it under the same balance band is the standard
approach taken by multilevel partitioners~\cite{metis,kahypar}, and Orion
implements it in three classical stages:
\begin{enumerate}
\item \emph{Initialization.} A geometric $k$-means clustering of the upper
  vectors into $P$ shards yields a cheap first assignment that ignores graph
  structure but is already roughly balanced.
\item \emph{Move-based refinement.} Following the Kernighan--Lin /
  Fiduccia--Mattheyses family~\cite{kl1970,fm1982}, Orion repeatedly considers
  moving a single upper point from its current shard to another. A move's
  \emph{gain} is the change in weighted cut (plus a topology vote from
  navigation evidence); a move is accepted only if its gain is nonnegative (or
  repairs a balance violation) and the target shard still has spare capacity.
  Sequential, capacity-aware acceptance avoids the classic collapse where many
  nodes race into the same underfull shard in one synchronous round.
\item \emph{Balance repair.} If some shards still violate the
  $(1{\pm}\varepsilon)$ band after refinement, Orion solves a small
  \emph{transportation} (min-cost flow) problem: overloaded shards are sources,
  underloaded shards are sinks, and admissible arcs are those justified by
  navigation evidence. Shipping one unit of flow along an arc corresponds to
  reassigning one point (or a group of points); the optimum restores balance at
  minimal disruption to the cut.
\end{enumerate}
All three stages read only the frozen upper graph, so the offline cost stays
near-linear and the result is checksum-reproducible.

\emph{Online}, serving is two-stage: a global traversal of $G$ selects shards and
per-shard entry points $\mathrm{EP}(q)\subseteq\mathrm{Vis}(q)$ (the estimator of
the neighbor-bearing shards), after which each selected shard runs an
\emph{independent} local search whose \emph{strength} scales with the same
evidence: a shard receiving more of $q$'s entry points is searched at a larger
$\mathrm{ef}_s=\mathrm{ef}_0+\gamma\,|\mathrm{EP}(q)\cap\pi^{-1}(s)|$. The
navigation signal thus governs both \emph{which} shards are probed (the fan-out)
and \emph{how hard} each is searched (the effort budget $\sum_s\mathrm{ef}_s$),
concentrating work on the shards navigation implicates. The \emph{realized}
fan-out is
$\mathrm{fanout}^\pi(q)=|\{\pi(v):v\in\mathrm{EP}(q)\}|$; base-layer placement
follows $\pi$ under the same balance band.

\paragraph{A mechanism-side bound and a recall assumption.}
Two claims connect Orion's mechanism back to the problem. First, the offline
objective it actually optimizes---the navigation-weighted cut---upper-bounds the
\emph{realized} online fan-out (Lemma~\ref{lem:fanout}). Second, minimizing that
realized fan-out is meaningful only if the router still returns the neighbors,
which we state as an explicit, empirically testable assumption
(Assumption~\ref{ass:recall}) rather than fold into the definitions.

\begin{lemma}[Realized fan-out is dominated by the weighted cut]
\label{lem:fanout}
Let $\mathrm{cut}(\pi)=\{(u,v)\in E:\pi(u)\neq\pi(v)\}$ and let $w$ be the
navigation traffic of~\eqref{eq:traffic}. Then
\[
\mathbb{E}_{q}\!\left[\mathrm{fanout}^\pi(q)\right]
\ \le\ 1+\!\!\sum_{(u,v)\in\mathrm{cut}(\pi)}\!\! w(u,v).
\]
\end{lemma}

\begin{proof}
Fix a query $q$, let $\mathrm{Vis}(q)$ be the set of nodes its upper traversal
expands, and let $T(q)\subseteq E$ be the set of edges the traversal explores.
Best-first search discovers a node only by exploring an edge out of an
already-expanded node, so $(\mathrm{Vis}(q),T(q))$ is \emph{connected}: it
contains a spanning tree rooted at the entry node. We stress that we do not
assume the expansion order forms a walk in $G$---consecutive expanded nodes need
not be adjacent, since candidates are popped from a priority queue.

Let $m=\bigl|\{\pi(v):v\in\mathrm{Vis}(q)\}\bigr|$ be the number of shards the
visited set spans, and contract each shard's portion of $\mathrm{Vis}(q)$ into a
single supernode. The contracted multigraph has $m$ supernodes and inherits
connectivity, hence carries at least $m-1$ edges; each of them is an explored
edge whose endpoints lie in different shards. Therefore
\[
m\ \le\ 1+\bigl|\{(u,v)\in T(q):\pi(u)\neq\pi(v)\}\bigr|.
\]
Entry points are selected only among expanded nodes, so
$\mathrm{EP}(q)\subseteq\mathrm{Vis}(q)$ and $\mathrm{fanout}^\pi(q)\le m$. Taking
expectations over $q$ and applying~\eqref{eq:traffic} edgewise,
$\mathbb{E}_q\bigl|\{(u,v)\in T(q):\pi(u)\neq\pi(v)\}\bigr|
=\sum_{(u,v)\in\mathrm{cut}(\pi)}w(u,v)$, which yields the claim.
\end{proof}

The bound is loose in two safe directions: fan-out counts only the shards of the
selected entry points, a subset of the visited shards, and a traversal that
repeatedly explores edges between the same two shards pays each time in the cut
but only once in fan-out. Thus the navigation-weighted cut is a
\emph{conservative} upper bound on realized fan-out, so optimizing it never
underestimates online cost.

A natural refinement of the same idea normalizes each shard's cut by its
\emph{volume} (total incident traffic)
$\mathrm{vol}(V_i)=\sum_{u\in V_i}\sum_{v}w(u,v)$, yielding the
\emph{conductance} (or normalized-cut) of the partition:
\[
\phi(\pi)=\sum_{i=1}^{P}\frac{\sum_{(u,v):\pi(u)=i,\pi(v)\neq i}w(u,v)}{\mathrm{vol}(V_i)}.
\]
Each summand is the probability that a random navigation step taken inside shard
$i$ immediately leaves it, and $\phi$ sums these escape probabilities over shards
(the normalized cut of Shi and Malik~\cite{shimalik2000}). Low conductance
therefore means every
shard is a \emph{metastable} region of the upper traversal---a set navigation
enters and then tends to stay in. Under this reading, the search-induced-topology hypothesis
simply says: shards that are metastable for navigation are the ones whose
interior points are co-relevant for the same queries, and are therefore the
right blocks for~\eqref{eq:objective}.

\begin{assumption}[Navigation preserves recall]
\label{ass:recall}
For target recall $R$, with high probability over $q\sim Q$ the shards selected by
the upper traversal, $\{\pi(v):v\in\mathrm{EP}(q)\}$, contain a recall-$R$
sufficient subset of $\mathrm{NN}_k(q)$.
\end{assumption}

Under Assumption~\ref{ass:recall}, Orion's router attains the target recall, so
its realized fan-out $\mathrm{fanout}^\pi$ is an \emph{achievable} cost for the
sharding problem, and by Lemma~\ref{lem:fanout} that cost is controlled by the
offline cut Orion minimizes. The required fan-out
$\mathrm{fanout}^\star_\sigma$ remains the mechanism-free reference against which
we report the residual gap; neither direction of comparison is automatic, since a
router targeting $R<1$ may deliberately skip a shard that holds a true neighbor.
The navigation graph thus links the mechanism-independent
objective to a computable surrogate \emph{without} being assumed in the problem
statement; Assumption~\ref{ass:recall} is what \S\ref{sec:eval} measures, and it
can fail---in which case the surrogate is simply a worse proxy, not a redefinition
of the problem.

\paragraph{From theory to mechanism.}
Table~\ref{tab:theory-map} maps each quantity in the formalization to the
concrete Orion mechanism and to the metric we report, so the design and the
evaluation can be read against the same objects.

One further design choice maps to a standard graph-partitioning relaxation.
So far $\sigma$ assigns each point to a single shard (\emph{edge-cut} style:
the cut is paid by edges that cross the boundary). Orion's
\emph{multi-assignment} instead lets a point occupy several shards, which is
the \emph{vertex-cut} (or overlapping) relaxation: a point that sits on a
search-relevant boundary is \emph{replicated} into the shards that navigation
evidence nominates. Replication lowers $\lambda$ for the hyperedges that
contain that point---the query no longer needs an extra hop to reach a copy---at
the cost of storing more physical copies. Formally, replication generalizes
$\sigma$ to $X\to 2^{\{1,\dots,P\}}$, the required fan-out becomes the smallest
number of shards \emph{covering} $\mathrm{NN}_k(q)$, and shard load is measured in
physical copies: with $M=\sum_{x\in X}|\sigma(x)|$ total copies, the balance band
in~\eqref{eq:objective} becomes $(1\pm\varepsilon)M/P$. The expansion ratio $M/N$
is the price of this relaxation, and is reported alongside fan-out so the
trade-off is explicit.

\begin{table}[t]
\centering
\small
\begin{tabular}{@{}lll@{}}
\toprule
Theoretical quantity & Orion mechanism & Reported metric \\
\midrule
Weighted cut / conductance & CCNB offline objective & cross-shard traffic \\
$\mathbb{E}[\mathrm{fanout}^\pi]$ (Lem.~\ref{lem:fanout}) & selective routing & visited logical shards \\
Balance band $(1{\pm}\varepsilon)$ & fixed-$P$ balancing & max/mean, CV \\
Load skew $\rho$ (2nd order) & size balance as proxy & max/mean query load \\
Per-shard strength $\mathrm{ef}_s$ (2nd order) & effort-scaled local search & within-shard work \\
Vertex-cut replication & multi-assignment & expansion ratio \\
\bottomrule
\end{tabular}
\caption{Correspondence between the partitioning formalization
(\S\ref{sec:theory}) and Orion's system mechanisms and measured quantities.}
\label{tab:theory-map}
\end{table}

We do not claim an approximation guarantee; instead, \S\ref{sec:eval}
empirically validates Lemma~\ref{lem:fanout} by correlating the offline
navigation-weighted cut (and conductance) with measured fan-out and end-to-end
QPS across $P\in\{4,8,16,32\}$, and tests Assumption~\ref{ass:recall} by checking
that the selected shards cover the true neighbors at the target recall. It also
quantifies the two second-order factors: the residual load skew $\rho$ that
survives size balancing, and the within-shard work that the effort-scaled
$\mathrm{ef}_s$ saves against a uniform-$\mathrm{ef}$ baseline at matched recall.

% ------------------------- STANDALONE (uncomment) --------------------------
% \end{document}
% ---------------------------------------------------------------------------
